\chapter{Thiết kế Khả kiểm bằng Phương pháp Quét (Design for Test by Means of Scan)}

\textit{Dịch và biên soạn chi tiết từ Chương 7 - ``Design for Test by Means of Scan'', kết hợp bổ sung phân tích chuyên sâu của nhóm.}

% ======================================================================
% MỤC 1: PHẦN BỔ SUNG CỦA NHÓM - KIẾN THỨC NỀN TẢNG
% ======================================================================
\section{Kiến thức nền tảng về Scan Design (Phần bổ sung của nhóm)}

\subsection{Tại sao cần DFT?}
Trong ngành công nghiệp bán dẫn hiện đại, mật độ tích hợp của các vi mạch (IC) đã đạt tới hàng tỷ transistor trên một chip duy nhất. Khi vi mạch được sản xuất xong, nhà sản xuất cần đảm bảo rằng mỗi chip đều hoạt động đúng chức năng trước khi giao cho khách hàng. Quá trình này gọi là \textbf{kiểm thử sản xuất (Manufacturing Test)}.

Đối với \textbf{mạch tổ hợp (Combinational Circuits)} thuần tuý, việc kiểm thử tương đối dễ dàng: ta chỉ cần đặt giá trị lên các ngõ vào (Primary Inputs) và đo kết quả tại các ngõ ra (Primary Outputs). Tuy nhiên, đối với \textbf{mạch tuần tự (Sequential Circuits)} -- loại mạch chiếm đa số trong các thiết kế thực tế -- thì bài toán phức tạp hơn rất nhiều. Nguyên nhân là bởi mạch tuần tự có chứa các \textbf{phần tử nhớ (Flip-Flops)} lưu giữ trạng thái bên trong. Các trạng thái này không thể được điều khiển trực tiếp từ bên ngoài, và cũng không thể quan sát trực tiếp tại các chân ra.

Kỹ thuật \textbf{Thiết kế Khả kiểm (Design for Test -- DFT)} ra đời nhằm giải quyết nút thắt này. Ý tưởng trung tâm của DFT là: \textit{chủ động sửa đổi cấu trúc phần cứng ngay từ giai đoạn thiết kế, để việc kiểm thử sau này trở nên dễ dàng và hiệu quả hơn}.

\subsection{Cấu trúc Scan Flip-Flop (MUX-DFF)}
Phương pháp phổ biến nhất trong DFT là \textbf{Scan Design}. Nguyên lý hoạt động của Scan Design dựa trên việc \textbf{thay thế} các D-Flip-Flop (DFF) thông thường trong mạch bằng các \textbf{Scan Flip-Flop (SFF)}, hay còn gọi là \textbf{MUX-DFF}.

\begin{figure}[!ht]
    \centering
    \includegraphics[width=0.55\textwidth]{images/mux_dff.png}
    \caption{Cấu trúc một Scan Flip-Flop (MUX-DFF) -- Phần bổ sung của nhóm}
    \label{fig:mux_dff}
\end{figure}

Như minh hoạ ở Hình \ref{fig:mux_dff}, một \textbf{bộ chọn kênh 2:1 (2-to-1 Multiplexer)} được gắn ngay trước ngõ vào D của Flip-Flop. Bộ chọn kênh này có chân điều khiển gọi là \texttt{SE} (Scan Enable) hoặc \texttt{NbarT} (Normal/Test mode). Tuỳ vào giá trị của tín hiệu điều khiển này, mạch sẽ hoạt động ở một trong hai chế độ:

\begin{itemize}
    \item \textbf{Chế độ hoạt động bình thường (Normal Mode, SE = 0):} Bộ chọn kênh cho phép tín hiệu từ đường dữ liệu logic (Data In) đi qua. Flip-Flop hoạt động như bình thường trong hệ thống, nhận dữ liệu từ mạch tổ hợp phía trước.
    \item \textbf{Chế độ kiểm thử (Test/Shift Mode, SE = 1):} Bộ chọn kênh \textit{ngắt} đường dữ liệu logic và thay vào đó cho phép tín hiệu từ chân \texttt{SI} (Scan Input) đi qua. Khi tất cả các SFF trong mạch đều ở chế độ này, chúng liên kết với nhau thành một \textbf{thanh ghi dịch (Shift Register)} dài. Ta có thể ``dịch'' dữ liệu vào (Shift In) qua chân \texttt{SI} và ``dịch'' dữ liệu ra (Shift Out) để quan sát.
\end{itemize}

Nhờ cơ chế này, mạch tuần tự phức tạp được \textbf{``tổ hợp hoá''} (combinationalized) -- tức là từ góc nhìn kiểm thử, các trạng thái nội tại trước đây ``ẩn'' bên trong giờ đây đều có thể được \textbf{điều khiển} (controllable) và \textbf{quan sát} (observable) thông qua đường scan.

\subsection{Quy trình kiểm thử Shift-Capture-Shift}
\textbf{[Bổ sung kiến thức của nhóm]} Toàn bộ quy trình kiểm thử bằng Scan tuân theo chu trình 3 bước lặp đi lặp lại cho mỗi test vector:

\begin{enumerate}
    \item \textbf{Shift-In:} Chuyển mạch sang chế độ Test (SE = 1). Dịch tuần tự từng bit của test vector vào chuỗi scan thông qua chân SI. Quá trình này tốn $L$ chu kỳ xung nhịp (clock cycles), trong đó $L$ là chiều dài chuỗi scan (số lượng SFF).
    \item \textbf{Capture:} Chuyển mạch sang chế độ Normal (SE = 0) trong đúng \textbf{1 chu kỳ clock}. Lúc này, mạch tổ hợp tính toán dựa trên giá trị vừa được nạp vào các Flip-Flop, và kết quả được ``chốt'' (capture) lại vào chính các Flip-Flop đó.
    \item \textbf{Shift-Out:} Chuyển lại sang chế độ Test (SE = 1). Dịch tuần tự các giá trị đã được capture ra ngoài qua chân SO để đem so sánh với kết quả mong đợi. Đồng thời, trong lúc shift-out, ta có thể shift-in test vector tiếp theo, giúp tiết kiệm thời gian.
\end{enumerate}

% ======================================================================
% MỤC 2: FULL SCAN TRÊN ADDING MACHINE (Dịch từ giáo trình)
% ======================================================================
\section{Kiểm thử Quét Toàn cục trên Adding Machine (7.5.1 Full Scan Adding Machine)}

\subsection{Cấu hình mạch sau khi chèn Scan}
Với cơ chế quét (scan) được chèn vào hệ thống Adding Machine như thể hiện ở Hình 7.52 trong sách gốc, đường dẫn tín hiệu \texttt{data\_bus\_in} trở thành ngõ vào chính (primary input) của mạch, còn \texttt{data\_bus\_out} cùng với \texttt{ad\_bus} hợp lại tạo thành cụm ngõ ra chính (primary output). Mạch khả kiểm thu được có tổng cộng \textbf{24 ngõ vào giả (pseudo primary inputs)} và \textbf{24 ngõ ra giả (pseudo primary outputs)}, lần lượt là các ngõ ra và ngõ vào của các thanh ghi (registers) nội tại.

\subsection{Sinh Test Vector bằng HOPE}
Để phục vụ kiểm thử mạch này, một \textbf{mô hình trải phẳng (unfolded model)} của toàn mạch đã được tạo ra, hoàn toàn tương tự cách đã thực hiện cho ví dụ mạch Residue-5 ở Mục 7.3 của sách. Mạng liên kết tổ hợp (combinational netlist) thu được sau khi trải phẳng được đưa vào công cụ sinh mẫu kiểm thử tự động (ATPG) có tên là \textbf{HOPE} để phân tích cho tổng cộng 976 lỗi (faults).

Kết quả mà HOPE trả về:
\begin{itemize}
    \item Tạo ra \textbf{64 test vectors} (bộ vector kiểm thử).
    \item Đạt \textbf{độ phủ lỗi (fault coverage) = 80.53\%}.
    \item Mỗi test vector gồm \textbf{32 bit ngõ vào} và \textbf{38 bit ngõ ra}:
    \begin{itemize}
        \item Phía input: 8 bit dành cho \texttt{data\_bus\_in} (đường song song) + 24 bit dành cho đường scan (đường nối tiếp).
        \item Phía output: 8 bit từ \texttt{data\_bus\_out} + 6 bit từ \texttt{ad\_bus} + 24 bit từ đường scan ra.
    \end{itemize}
\end{itemize}

\subsection{Mô hình Tester ảo (Virtual Tester)}
Mô hình khả kiểm hoàn chỉnh của Adding Machine thu được bằng cách xâu chuỗi tất cả các flip-flop của mạch thành \textbf{một chuỗi scan duy nhất (single chain)}. Hình \ref{lst:753} trình bày bộ tester ảo (virtual tester) dùng cho việc kiểm thử full scan trên mạch khả kiểm này. Ngoại trừ mạch đang kiểm thử (circuit under test) là khác biệt, và số lần dịch (shifts) trong các scan flip-flop khác nhau, \textbf{bộ tester này hoàn toàn tương tự với bộ tester của mạch Residue-5} đã được thảo luận ở Mục 7.3 của sách. Hình \ref{lst:753} chỉ hiển thị phần vòng lặp chịu trách nhiệm áp dụng các test vector vào netlist đã chèn scan.

\subsection{Giải thích chi tiết Testbench}
Testbench này đọc lần lượt \textbf{64 test vectors} từ tệp \texttt{testFile} và áp dụng chúng vào mạch full scan để kiểm tra từng lỗi trong số 976 lỗi. Với mỗi test vector:
\begin{enumerate}
    \item \textbf{8 bit} được áp vào đường song song (\texttt{PI}).
    \item \textbf{24 bit còn lại} được dịch vào nối tiếp (shifted-in serially) qua vòng lặp \texttt{repeat(24)} -- mỗi lần lặp đẩy 1 bit vào chân \texttt{si} và đồng thời rút 1 bit từ chân \texttt{so}.
    \item Sau khi dịch hết 24 bit, mạch chuyển về chế độ Normal (\texttt{NbarT = 0}) trong \textbf{1 xung clock} để thực hiện bước Capture.
    \item Khối lệnh \texttt{if} so sánh chuỗi ghép nối (concatenation) giữa 24 bit output thu thập từ đường nối tiếp và 14 bit output song song với 38 bit đầu ra mong đợi từ \texttt{testFile}. Nếu khác nhau $\Rightarrow$ lỗi đã được phát hiện (detected).
\end{enumerate}

\begin{lstlisting}[caption={Hình 7.53: Bộ tester ảo cho Full Scan Adding Machine (Full scan Adding Machine virtual tester)}, label={lst:753}]
module Tester;
. . .. . .
CPU_ScanInserted FUT(clk, PI, PO, NbarT, Si, So);
always #200 clk = ~clk;
initial begin
    while( !$feof(faultFile)) begin
        . . .
        . . .
        while((!$feof(testFile))&&(!detected))begin
            status = $fscanf(testFile,"%b\n",line);
            pre_expected_st = cur_expected_st;
            expected_PO = line[out_size - 1:0];
            cur_expected_st = line[st_size + out_size -1 :out_size];
            PI = line[st_size+out_size+in_size-1:st_size+out_size];
            NbarT = 1'b1;
            #delay;
            index = stIndex;
            repeat(24) begin
                si = line[index];
                @(posedge clk);
                saved_st[index - stIndex] = so;
                index = index + 1;
            end
            NbarT = 1'b0;
            @(posedge clk);
            SampledPO = PO;
            if({pre_expected_st, expected_PO} != {saved_st, SampledPO})
            begin
                numOfDetected = numOfDetected + 1;
                detected = 1;
            end
            #5;
        end//test
        . . .
    end//fault
    . . .
end//initial
endmodule
\end{lstlisting}

\subsection{Phân tích nút thắt thời gian test (Bổ sung của nhóm)}
\textbf{[Bổ sung kiến thức của nhóm]} Vòng lặp \texttt{repeat(24)} chính là ``nút thắt cổ chai'' (bottleneck) lớn nhất của Full Scan. Với ví dụ Adding Machine chỉ có 24 flip-flop, thì thời gian dịch không đáng kể. Nhưng hãy xem xét công thức tổng quát cho thời gian test:

$$T_{test} \approx V \times (L + 1) \times T_{clk}$$

Trong đó: $V$ = số lượng test vectors, $L$ = chiều dài chuỗi scan (số flip-flop), $T_{clk}$ = chu kỳ xung nhịp.

Trong thiết kế chip thương mại hiện đại, một bộ xử lý có thể chứa hàng triệu flip-flop. Nếu áp dụng Full Scan dạng chuỗi đơn (single chain), giả sử $L = 1.000.000$ Flip-flop, $V = 10.000$ vectors, $T_{clk} = 10 ns$:

$$T_{test} = 10.000 \times 1.000.001 \times 10 \text{ ns} = 100.000 \text{ giây} \approx 27.8 \text{ giờ}$$

Con số này là không thể chấp nhận được trong sản xuất hàng loạt, nơi mỗi giây trên máy ATE (Automatic Test Equipment) đều tốn chi phí rất lớn. Đây chính là lý do thúc đẩy sự ra đời của kỹ thuật \textbf{Multiple Parallel Scans} ở mục tiếp theo.

% ======================================================================
% MỤC 3: MULTIPLE SCAN TRÊN ADDING MACHINE (Dịch từ giáo trình)
% ======================================================================
\section{Thiết kế Quét Đa luồng cấp RTL (7.5.2 RTL Design Multiple Scan)}

\subsection{Ý tưởng phân chia chuỗi}
Mục này trình bày cách triển khai thiết kế quét đa luồng (multiple scan design) trên mạch Adding Machine đã giới thiệu ở Chương 2 của sách gốc. Tương tự phần thảo luận về Full Scan ở mục trước, mặc dù sơ đồ mạch trông khác so với mô hình Huffman ở Hình 7.44, nhưng bản chất cấu trúc mạch vẫn tuân theo mô hình đó, và phương pháp quét đa luồng (multiple scan) ở Mục 7.4 hoàn toàn áp dụng được cho mạch này.

\subsection{Thiết kế Quét ở Mức RT (7.5 RT Level Scan Design)}
Chúng ta sử dụng \textbf{3 chuỗi scan có chiều dài bằng nhau}:
\begin{enumerate}
    \item \textbf{Chuỗi 1:} Bao phủ thanh ghi tích luỹ \texttt{AC} (Accumulator).
    \item \textbf{Chuỗi 2:} Bao phủ thanh ghi lệnh \texttt{IR} (Instruction Register).
    \item \textbf{Chuỗi 3:} Bao phủ thanh ghi bộ đếm chương trình \texttt{PC} (Program Counter) cộng thêm 2 flip-flop điều khiển.
\end{enumerate}

Vì kích thước cả 3 chuỗi scan bằng nhau, chúng ta có thể áp dụng cấu hình \textbf{quét song song đa chuỗi (multiple parallel scan chains)}, giúp đơn giản hoá đáng kể logic điều khiển scan. Hình \ref{fig:scan_chain} minh hoạ 3 chuỗi scan này trong sơ đồ RTL của Adding Machine.

\begin{figure}[!ht]
    \centering
    \includegraphics[width=0.75\textwidth]{images/scan_chain.png}
    \caption{Hình 7.54: Ba chuỗi scan song song cho Adding Machine (Multiple parallel scans for Adding Machine)}
    \label{fig:scan_chain}
\end{figure}

\subsection{Chèn chuỗi Scan vào Netlist}
Các chuỗi scan được \textbf{chèn thủ công (manually inserted)} vào netlist. Netlist này chính là kết quả của quá trình tổng hợp (synthesize) mã Verilog của Adding Machine. Việc chèn scan sẽ gây ra thay đổi ở \textbf{hai chỗ}:
\begin{enumerate}
    \item \textbf{Cổng giao tiếp (ports):} Thêm các chân vào/ra mới cho đường scan.
    \item \textbf{Sắp xếp flip-flop:} Các flip-flop được kết nối lại thành các chuỗi theo thứ tự mong muốn.
\end{enumerate}

Hình \ref{lst:755} hiển thị một phần mã nguồn của netlist đã được sửa đổi. Cụ thể:
\begin{itemize}
    \item \textbf{Ngõ vào mới:} Chân \texttt{NbarT} (chuyển chế độ Normal/Test) và 3 chân serial input: \texttt{ir\_Si}, \texttt{ac\_Si}, \texttt{pc\_Si}.
    \item \textbf{Ngõ ra mới:} 3 chân serial scan output: \texttt{ir\_So}, \texttt{ac\_So}, \texttt{cntrl\_So}.
    \item \textbf{3 khu vực được đánh dấu bằng dấu //} trong mã nguồn tương ứng với 3 chuỗi scan:
    \begin{itemize}
        \item Chuỗi 1: Bắt đầu từ \texttt{ir\_Si}, đi qua \texttt{INS\_1} đến \texttt{INS\_8}, kết thúc tại \texttt{wire\_26} $\rightarrow$ gán cho \texttt{ir\_So}.
        \item Chuỗi 2: Bắt đầu từ \texttt{ac\_Si}, đi qua \texttt{INS\_9} đến \texttt{INS\_16}, kết thúc tại \texttt{wire\_52} $\rightarrow$ gán cho \texttt{ac\_So}.
        \item Chuỗi 3: Bắt đầu từ \texttt{pc\_Si}, đi qua \texttt{INS\_17} đến \texttt{INS\_24}, kết thúc tại \texttt{wire\_4} $\rightarrow$ gán cho \texttt{cntrl\_So}.
    \end{itemize}
    \item \textbf{Tất cả} các flip-flop đều sử dụng \texttt{NbarT} làm tín hiệu đầu vào điều khiển chế độ.
\end{itemize}

\begin{lstlisting}[caption={Hình 7.55: Chèn ba chuỗi quét vào netlist (Inserting three scan chains)}, label={lst:755}]
module CPU_M_ScanInserted(clk, {reset, data_bus_in}, {adr_bus,rd_mem,
wr_mem,data_bus_out}, NbarT, ir_Si, ac_Si,
pc_Si, ir_So, ac_So, cntrl_So);
. . .
. . .
. . .
. . .
dff INS_1 (wire_6,wire_96,clk,1'b0,1'b0,ir_en,NbarT,ir_Si, 1'b0);
dff INS_2 (wire_2,wire_98,clk,1'b0,1'b0,ir_en,NbarT,wire_6, 1'b0);
dff INS_8 (wire_26,wire_173,clk,1'b0,1'b0,ir_en,NbarT,wire_54, 1'b0);
//
dff INS_9 (wire_130,wire_180,clk,1'b0,1'b0,ac_en,NbarT,ac_Si, 1'b0);
dff INS_10(wire_135,wire_186,clk,1'b0,1'b0,ac_en,NbarT,wire_130, 1'b0);
dff INS_16(wire_52,wire_208,clk,1'b0,1'b0,ac_en,NbarT,wire_50, 1'b0);
//
dff INS_17 (wire_65,wire_211,clk,1'b0,1'b0,pc_en,NbarT,pc_Si, 1'b0);
dff INS_18(wire_71,wire_216,clk,1'b0,1'b0,pc_en,NbarT,wire_65, 1'b0);
dff INS_24 (wire_4,wire_249,clk,1'b0,1'b0,pc_en,NbarT,wire_18, 1'b0);

assign ir_So = wire_26;
assign ac_So = wire_52;
assign cntrl_So = wire_4;
endmodule
\end{lstlisting}

\subsection{Chương trình Test cho Multiple Scan}
Để viết chương trình test cho thiết kế này, ta cần lưu ý rằng thiết kế có \textbf{3 chuỗi scan có cùng chiều dài}, sử dụng \textbf{chung một xung nhịp (clock)}, và \textbf{chung một tín hiệu điều khiển chế độ} \texttt{NbarT}.

Hình \ref{lst:756} trình bày bộ tester ảo (virtual tester) mô phỏng hành vi của một thiết bị ATE và chương trình test tương ứng. Hình này chỉ hiển thị phần khởi tạo (instantiation) mạch 3 chuỗi scan và thủ tục áp dụng dữ liệu test. Các phần khác của testbench về cơ bản giống hệt với mã đã thảo luận ở Mục 7.3, nên không lặp lại ở đây.

Như thể hiện trong mã, khối instantiation của netlist (mã một phần hiển thị ở Hình \ref{lst:755}) bao gồm \texttt{NbarT} cùng các ngõ vào và ngõ ra nối tiếp của cả 3 chuỗi scan.

\subsection{Giải thích chi tiết quy trình Test}
Thủ tục test được đặt bên trong khối \texttt{initial} của Hình \ref{lst:756}. Tất cả các phần đều giống hệt Hình 7.53 (\ref{lst:753}), \textbf{ngoại trừ vòng lặp \texttt{repeat}}. Vòng lặp \texttt{repeat} ở đây hoạt động như sau:

\begin{enumerate}
    \item Lấy từng nhóm \textbf{8 bit} từ biến \texttt{line} (chứa test vector đầu vào) và gán đồng thời cho 3 chân scan input: \texttt{ir\_Si}, \texttt{ac\_Si}, \texttt{pc\_Si}.
    \item Việc gán này diễn ra khi \texttt{NbarT = 1} cho cả 3 chuỗi (tức là ở chế độ Shift).
    \item Sau mỗi lần gán, các thanh ghi (registers) được cấp xung clock \textbf{đồng thời (simultaneously)}.
    \item Sau xung clock, kết quả đầu ra từ \texttt{ir\_So}, \texttt{ac\_So}, \texttt{cntrl\_So} được thu thập và lưu vào mảng \texttt{saved\_st}.
    \item Các giá trị output được \textbf{xếp cách nhau 8 bit} trong mảng \texttt{saved\_st} để đảm bảo các bit đầu ra của cùng một chuỗi scan nằm liền kề nhau.
\end{enumerate}

\begin{lstlisting}[caption={Hình 7.56: Chương trình test cho quét đa luồng song song (Test program for multiple parallel scans)}, label={lst:756}]
module Tester;
. . .
. . .
. . .
CPU_M_ScanInserted FUT(clk, PI, PO, NbarT, ir_Si, ac_Si, pc_Si, ir_So,
ac_So, cntrl_So);
always #200 clk = ~clk;
initial begin
    while( !$feof(faultFile)) begin
        while((!$feof(testFile))&&(!detected)) begin
            status = $fscanf(testFile,"%b\n",line);
            pre_expected_st = cur_expected_st;
            expected_PO = line[out_size - 1:0];
            cur_expected_st = line[st_size + out_size -1 :out_size];
            PI = line[st_size+out_size+in_size-1:st_size+out_size];
            NbarT = 1'b1;
            #delay;
            index = stIndex;
            repeat(8) begin
                ir_Si = line[index];
                ac_Si = line[index+8];
                pc_Si = line[index+16];
                @(posedge clk);
                saved_st[index - stIndex] = ir_So;
                saved_st[index+8 - stIndex] = ac_So;
                saved_st[index+16 - stIndex] = cntrl_So;
                index = index + 1;
            end
            NbarT = 1'b0;
            @(posedge clk);
            SampledPO = PO;
            if({pre_expected_st, expected_PO} != {saved_st, SampledPO})
            begin
                numOfDetected = numOfDetected + 1;
                detected = 1;
            end
            #5;
        end//test
        . . .
    end//fault
    . . .
end//initial
endmodule
\end{lstlisting}

\subsection{So sánh hiệu quả Full Scan vs Multiple Scan (Bổ sung của nhóm)}
\textbf{[Bổ sung kiến thức của nhóm]} Bảng dưới đây tổng hợp sự khác biệt then chốt giữa hai phương pháp khi áp dụng cho mạch Adding Machine (24 flip-flop, 64 test vectors):

\begin{table}[!ht]
\centering
\caption{So sánh Full Scan và Multiple Scan trên Adding Machine}
\label{tab:compare}
\begin{tabular}{|l|c|c|}
\hline
\textbf{Tiêu chí} & \textbf{Full Scan (1 chain)} & \textbf{Multiple Scan (3 chains)} \\
\hline
Số chuỗi scan & 1 & 3 \\
\hline
Chiều dài mỗi chuỗi & 24 flip-flop & 8 flip-flop \\
\hline
Số chu kỳ shift/vector & 24 & 8 \\
\hline
Tổng shift cycles (64 vectors) & $64 \times 24 = 1536$ & $64 \times 8 = 512$ \\
\hline
Số chân Scan In & 1 & 3 \\
\hline
Số chân Scan Out & 1 & 3 \\
\hline
Tổng chân I/O tốn thêm & 2 & 6 \\
\hline
Tốc độ test & Chậm ($\times 1$) & Nhanh gấp 3 ($\times 3$) \\
\hline
\end{tabular}
\end{table}

Công thức tổng quát khi chia thành $N$ chuỗi song song bằng nhau:
$$T_{multiple} = V \times \frac{L}{N} \times T_{clk}$$
So với Full Scan: $T_{full} = V \times L \times T_{clk}$. Hệ số giảm thời gian là $N$ lần, nhưng phải \textbf{đánh đổi bằng $2 \times N$ chân I/O bổ sung} (mỗi chuỗi cần 1 SI và 1 SO).

% ======================================================================
% MỤC 4: CÁC MÔ HÌNH SCAN CẤP RTL (Dịch từ giáo trình)
% ======================================================================
\section{Quy hoạch các mô hình Scan cho thiết kế cấp RTL (7.5.3 Scan Designs for RTL)}

Việc lập kế hoạch (planning) cho các thiết kế scan khác nhau ở \textbf{mức RT-level} là tương đối dễ dàng nếu ta \textbf{nhìn nhận mạch từ góc độ RTL view} -- tức là xem mạch dưới dạng các thanh ghi (registers) và các khối logic tổ hợp kết nối chúng. Khi đã lập kế hoạch xong, ta có thể \textbf{ước lượng thời gian test} (test time estimate). Việc triển khai thực tế (implementation) thiết kế scan đã lên kế hoạch sẽ được thực hiện trên \textbf{netlist đầu ra} của quá trình tổng hợp (synthesis) từ mã HDL cấp RT-level. Việc phát triển chương trình test cũng trở nên dễ dàng hơn khi ta suy nghĩ về mạch dưới dạng \textbf{mô hình Huffman} sau khi đã chèn scan.

\subsection{Partial Scan (Quét Một phần)}
Ngoài Full Scan và Multiple Scan, còn có các phương án thiết kế scan khác cho ví dụ RTL của chúng ta, đó là \textbf{Partial Scan} (quét một phần). Trong Partial Scan, \textbf{không phải tất cả} flip-flop đều được thay bằng Scan Flip-Flop. Chỉ một số flip-flop được chọn lọc (thường là những flip-flop nằm trên các vòng lặp phản hồi -- feedback loops -- khó kiểm thử nhất) mới được chèn scan.

\textbf{[Bổ sung kiến thức của nhóm]} Ưu điểm của Partial Scan là giảm đáng kể diện tích phần cứng bổ sung (area overhead) và giảm ảnh hưởng đến hiệu năng (timing/power). Tuy nhiên, nhược điểm là \textbf{độ phủ lỗi (fault coverage)} thường thấp hơn Full Scan, và công cụ ATPG phải xử lý mô hình phức tạp hơn (vì mô hình không hoàn toàn là tổ hợp -- vẫn còn một số thanh ghi).

\subsection{Multiple Independent Scans (Quét Đa luồng Độc lập)}
Một giải pháp khác là \textbf{Multiple Independent Scans} (quét đa luồng độc lập). Ý tưởng ở đây là nếu ta \textbf{phân chia (partition)} mạch thành hai phần riêng biệt -- ví dụ phần \texttt{Datapath} (đường dẫn dữ liệu) và phần \texttt{Controller} (bộ điều khiển) -- thì có thể thiết kế scan \textbf{riêng biệt, độc lập} cho từng phần.

\subsection{Isolated Scan (Quét Cô lập)}
Việc tách một mạch thành nhiều phần con (subcomponents) có thể được thực hiện thông qua kỹ thuật \textbf{Isolated Scan} (quét cô lập). Trong kỹ thuật này, một thanh ghi dịch (shift register) được đặt tại \textbf{giao diện (interface)} giữa các phần con, cho phép dịch vào (scan-in) và dịch ra (scan-out) các giá trị tín hiệu ở ranh giới giữa chúng. Nhờ vậy, mỗi phần con có thể được kiểm thử \textbf{cô lập}, độc lập với các phần khác.

\subsection{Các tham số quyết định khi lựa chọn chiến lược Scan}
Khi lập kế hoạch cho các chiến lược scan khác nhau, có \textbf{3 tham số chính} cần cân nhắc:
\begin{enumerate}
    \item \textbf{Số lượng test vectors:} Càng nhiều vector thì thời gian test càng lâu, nhưng độ phủ lỗi có thể cao hơn.
    \item \textbf{Chiều dài mỗi test vector:} Phụ thuộc vào chiều dài chuỗi scan và số lượng ngõ vào/ra song song.
    \item \textbf{Độ phủ lỗi (Fault Coverage):} Chỉ số quan trọng nhất, thể hiện tỷ lệ phần trăm lỗi có thể được phát hiện.
\end{enumerate}

\textbf{[Bổ sung kiến thức của nhóm]} Bảng dưới đây tóm tắt ưu nhược điểm của từng chiến lược:

\begin{table}[!ht]
\centering
\caption{So sánh các chiến lược Scan Design}
\label{tab:strategies}
\begin{tabular}{|l|c|c|c|}
\hline
\textbf{Chiến lược} & \textbf{Fault Coverage} & \textbf{Area Overhead} & \textbf{Test Time} \\
\hline
Full Scan (1 chain) & Cao nhất & Cao nhất & Rất lâu \\
\hline
Multiple Parallel Scan & Cao nhất & Cao (thêm chân I/O) & Giảm $N$ lần \\
\hline
Partial Scan & Trung bình & Thấp & Trung bình \\
\hline
Isolated Scan & Tuỳ thiết kế & Trung bình & Tuỳ phân vùng \\
\hline
\end{tabular}
\end{table}

% ======================================================================
% MỤC 5: TỔNG KẾT CHƯƠNG (Dịch từ giáo trình)
% ======================================================================
\section{Tổng kết chương (7.6 Summary)}

Chương này đã thảo luận về các \textbf{kỹ thuật DFT (Design for Test)}. Chúng ta đã bắt đầu với một số \textbf{phương pháp đặc thù (ad hoc methods)} và từ đó phát triển dần lên tới \textbf{Isolated Scan}. Tất cả các phương pháp này đều có thể áp dụng cho cả mạch tổ hợp (combinational) lẫn mạch tuần tự (sequential).

Kỹ thuật \textbf{Scan Design} mà chúng ta bắt đầu thảo luận từ Mục 7.3 chủ yếu tập trung vào \textbf{mạch tuần tự}, và mục đích chính của nó là \textbf{biến đổi mạch tuần tự thành mạch tổ hợp} (từ góc nhìn kiểm thử), sao cho các phương pháp sinh test tổ hợp (combinational TG methods) có thể được sử dụng để tạo mẫu kiểm thử.

Các thiết kế scan khác mà chúng ta đã thảo luận về cơ bản đều tuân theo \textbf{cùng một nguyên tắc}. Trong \textbf{Partial Scan methods}, việc đưa vào các ngõ vào giả (pseudo inputs) và ngõ ra giả (pseudo outputs) \textbf{không hoàn toàn} biến đổi mô hình tuần tự thành mô hình tổ hợp -- các thanh ghi vẫn còn tồn tại trong mô hình thu được. Mặc dù vậy, mô hình thu được \textbf{vẫn có thể} được sử dụng cho quá trình sinh test tổ hợp. Sự sai lệch (deviation) so với mô hình tổ hợp thuần tuý sẽ ảnh hưởng đến \textbf{thủ tục test} (test procedure), trong đó cần bổ sung thêm các \textbf{xung nhịp phụ (extra clock pulses)} để bù đắp cho những thanh ghi vẫn còn lại trong mô hình.

Để minh hoạ việc ứng dụng phương pháp scan cho \textbf{thiết kế cấp RT-level} và chứng minh tính khả thi của các kỹ thuật này, mục cuối cùng đã sử dụng một thiết kế RT-level nhỏ gọn kết hợp với \textbf{Verilog testbench} để kiểm thử scan đã chèn. Các thiết kế lớn hơn có thể được \textbf{phân chia (partition)} thành nhiều phần con (subcomponents), trong đó mỗi phần con có thể yêu cầu một \textbf{dạng scan khác nhau}. Việc tách biệt thiết kế thành các phần con có thể được thực hiện bằng \textbf{thanh ghi dịch (shift register)} để scan-in và scan-out các giá trị tín hiệu tại giao diện (interface) giữa các phần con.

% ======================================================================
% TÀI LIỆU THAM KHẢO
% ======================================================================
\section*{Tài liệu tham khảo (References)}
\begin{enumerate}
    \item Abramovici M, Breuer MA, Friedman AD (1994) \textit{Digital systems testing and testable design}. IEEE Press, Piscataway, NJ, revised printing.
    \item Wilkins BR (1986) \textit{Testing digital circuits, an introduction}. Van Nostrand Reinhold, Berkshire, UK.
    \item Eichelberger EB, Lindbloom E, Waicukauski JA, Williams TW (1991) \textit{Structured logic testing}. Prentice-Hall, Englewood Cliffs, NJ.
    \item Agrawal VD, Mercer MR (1982) Testability measures -- What do they tell us? In: \textit{Proceedings of the International Test Conference}, Nov 1982, pp 391--396.
    \item Williams MJY, Angell JB (1973) Enhancing testability of large-scale integrated circuits via test points and additional logic. \textit{IEEE Trans Comput} C-22(1):46--60.
    \item Miczo A (2003) \textit{Digital logic testing and simulation}, 2nd edn. Wiley, Hoboken, NJ.
    \item Cheng K-T, Lin C-J (1995) Timing-driven test point insertion for full-scan and partial-scan BIST. In: \textit{Proceedings of the International Test Conference}, Oct 1995, pp 506--514.
    \item Cheng K-T, Agrawal VD (1990) A partial scan method for sequential circuits with feedback. \textit{IEEE Trans Comput} 39(4):544--548.
    \item Narayanan S, Gupta R, Breuer MA (1993) Optimal configuring of multiple scan chains. \textit{IEEE Trans Comput} 42(9):1121--1131.
    \item Jha NK, Gupta S (2003) \textit{Testing of digital systems}. Cambridge University Press, Cambridge, UK.
\end{enumerate}
